% Experiment report. Fixed sections: see experiments/TEMPLATE/report.tex and experiments/README.md.
% Build: npm run experiment:report -- 19-hill-mirror-credit
\documentclass{../polyexp}
\begin{document}
% generated by scripts/experiment-report.ts from experiment.json and models/<run>/ - do not edit
\providecommand{\ExpHeader}{}\renewcommand{\ExpHeader}{%
{\sffamily\footnotesize\color{expgrey}EXPERIMENT 19 \textperiodcentered\ METHOD \hfill 2026-11-24}\par
{\LARGE\bfseries\sffamily Hill climbing that learns from lost rounds: mirror and credit}\hfill\verdictbox{expred}{NOT KEPT}\par\smallskip
{\itshape Does hill climbing find more improvements when it learns from the rounds it loses: trying the opposite direction after a clear loss (mirror), and choosing which knobs to nudge from a per-knob model of all past rounds (credit)?}\par\medskip
\begin{tabular}{@{}p{22mm}p{\dimexpr\linewidth-26mm}@{}}
\fact{Change}{--mirror: after a round lost by at least 1 point, the next round tests the same knob moves in the opposite direction. --credit: a per-knob linear model of the margin, fitted on every round and weighted by games, picks knobs and directions by Thompson sampling; 25\% of rounds stay random; the evidence is halved after each keep. Both off by default.}
\fact{Base}{\model{greedy-full-v1}}
\fact{Method}{Tuning run (hill climbing), 100,000 games, against a control run of the same budget and seed; metric: points share}
\fact{Setups}{4 players, 16$\times$16, drylands, turn cap 150}
\fact{Record}{issue \issue{31}; status: done}
\end{tabular}\par\medskip
\begin{center}\begin{tikzpicture}[node distance=5mm]
    \node[flowbox] (n0) {{\sffamily\bfseries\scriptsize\color{expgrey}BASE}\\[1pt]\model{full-v1}};
    \node[flowbox,right=5mm of n0] (n1) {{\sffamily\bfseries\scriptsize\color{expgrey}CHANGE}\\[1pt]method};
    \node[flowbox,right=5mm of n1] (n2) {{\sffamily\bfseries\scriptsize\color{expgrey}MEASURED}\\[1pt]Tuning run, 100,000 games, vs a control\\{\scriptsize 1 setup, metric: points share}};
    \node[flowbox,right=5mm of n2] (n3) {{\sffamily\bfseries\scriptsize\color{expgrey}RESULT}\\[1pt]Points share, 6-model panel, at 100,000 games\\{\scriptsize 66.1 $\rightarrow$ 66.6\,\%}\\\textbf{+0.5\pp}};
    \node[flowbox,right=5mm of n3] (n4) {{\sffamily\bfseries\scriptsize\color{expgrey}VERDICT}\\[1pt]\textbf{\color{expred}NOT KEPT}};
    \draw[flowarrow] (n0) -- (n1); \draw[flowarrow] (n1) -- (n2); \draw[flowarrow] (n2) -- (n3); \draw[flowarrow] (n3) -- (n4);
\end{tikzpicture}\end{center}
\begin{center}\sffamily\small\begin{tabular}{@{}lrrr@{}}\toprule
Measure & Base & New & Change\\\midrule
Points share, 6-model panel, at 100,000 games & 66.1\,\% & 66.6\,\% & +0.5\\
Points share, 6-model panel, start (greedy-full-v1) & 67.7\,\% & 67.7\,\% & +0.0\\
Changes kept in 100,000 games & 3.0 & 2.0 & $-$1.0\\
\bottomrule\end{tabular}\end{center}
}
\providecommand{\ExpResults}{}\renewcommand{\ExpResults}{%
\subsection{Mirror + credit, 100,000 games {\normalfont\small(\model{greedy-full-v2})}}
\begin{center}\begin{tikzpicture}\begin{axis}[expplot,height=4.6cm,xlabel={training games (thousands)},ylabel={Points share, panel mean (\%)},enlarge y limits=0.2]
\addplot[dashed,expgrey,line width=0.7pt,forget plot] coordinates {(0.0,67.70) (100.0,67.70)};
\addplot[expblue,mark=*,mark size=1.8pt] coordinates {(0.0,67.70) (12.9,66.30) (25.1,66.30) (37.8,66.30) (50.4,66.30) (62.6,66.30) (75.0,66.30) (87.9,66.60) (100.0,66.60)};
\end{axis}\end{tikzpicture}\\[-2pt]{\footnotesize\color{expgrey}Learning curve of the tuning run (2 kept changes in 315 rounds). Dashed: the start.}\end{center}
\begin{center}\sffamily\small\begin{tabular}{@{}lrr@{}}\toprule
Opponent & Won (\%) & Points (\%)\\\midrule
\model{v1} & 95 $\pm$ 2 & 75.0\\
\model{hill-v1} & 85 $\pm$ 3 & 68.6\\
\model{hill-v2} & 94 $\pm$ 2 & 71.4\\
\model{hill-v3} & 88 $\pm$ 3 & 67.9\\
\model{hill-v4} & 88 $\pm$ 3 & 67.2\\
\model{full-v1} & 50 $\pm$ 5 & 50.0\\
\midrule
\textbf{Mean} & 83.1 & 66.7\\
\bottomrule\end{tabular}\\[3pt]{\footnotesize\color{expgrey}Final test of the saved model: 400 unseen games against each panel model.}\end{center}

\subsection{Control (standard hill climbing), 100,000 games {\normalfont\small(\model{greedy-full-v2-ctrl})}}
\begin{center}\begin{tikzpicture}\begin{axis}[expplot,height=4.6cm,xlabel={training games (thousands)},ylabel={Points share, panel mean (\%)},enlarge y limits=0.2]
\addplot[dashed,expgrey,line width=0.7pt,forget plot] coordinates {(0.0,67.70) (100.4,67.70)};
\addplot[expblue,mark=*,mark size=1.8pt] coordinates {(0.0,67.70) (12.8,66.30) (25.1,66.30) (37.6,66.30) (50.4,66.30) (62.9,66.30) (75.3,66.30) (87.9,66.30) (100.4,66.10)};
\end{axis}\end{tikzpicture}\\[-2pt]{\footnotesize\color{expgrey}Learning curve of the tuning run (3 kept changes in 318 rounds). Dashed: the start.}\end{center}
\begin{center}\sffamily\small\begin{tabular}{@{}lrr@{}}\toprule
Opponent & Won (\%) & Points (\%)\\\midrule
\model{v1} & 95 $\pm$ 2 & 75.0\\
\model{hill-v1} & 85 $\pm$ 3 & 68.6\\
\model{hill-v2} & 94 $\pm$ 2 & 71.4\\
\model{hill-v3} & 88 $\pm$ 3 & 67.9\\
\model{hill-v4} & 88 $\pm$ 3 & 67.2\\
\model{full-v1} & 50 $\pm$ 5 & 50.0\\
\midrule
\textbf{Mean} & 83.1 & 66.7\\
\bottomrule\end{tabular}\\[3pt]{\footnotesize\color{expgrey}Final test of the saved model: 400 unseen games against each panel model.}\end{center}
}

\ExpHeader

\section{Question and hypothesis}
Run on 24 November, before experiment 1, and written up on 26 November as a record: no hypothesis was written down before the runs, so this section states the question only, not a prediction made in advance. The number is 19 because it took a number when it was recorded, not when it ran.

Hill climbing throws away what a lost round teaches: a challenger that loses clearly is dropped, and the next round draws new random knob moves. The question was whether a climber that learns from those rounds keeps more improvements in the same budget: by testing the opposite direction after a clear loss (\emph{mirror}), and by steering which knobs to move, and which way, from a model of all past rounds (\emph{credit}).
\todo{jakob}{The idea in your words, if you remember what you expected.}

\section{Mechanism}
Both are options of \texttt{scripts/train-hill.ts}, off by default (runs without them behave as before).
\texttt{--mirror}: after a round whose mean margin is at most \texttt{--mirror-below} (default $-1$ point), the next round tests the same knob moves with the sign flipped.
\texttt{--credit}: a per-knob linear model of the round margin, fitted on every round (kept or not) and weighted by its games, draws knobs and directions by Thompson sampling (prior sd 0.1); a share \texttt{--credit-explore} of rounds (25\,\%) stays random, and the evidence is halved after each keep.
Every round is also written to \texttt{models/<id>/rounds.jsonl}: its source (random, mirror, credit), the knobs moved, games, margin, log-likelihood ratio and whether it was kept.

\section{Setup}
Two tuning runs from \model{greedy-full-v1} with identical settings (100,000 games, seed 4242, league \texttt{self,past}, batch 30, curve every 12,500 games, panel of six), one with \texttt{--mirror --credit} (\model{greedy-full-v2}) and one without (\model{greedy-full-v2-ctrl}, the control that experiment 1 also uses). Both started at the same time on the same machine.

\section{Results}
\ExpResults

Both learning curves stay at or below the start: the panel points share was 67.7\,\% at 0 games and 66.6\,\% (mirror + credit) and 66.1\,\% (control) at 100,000 games; with the champion's weights averaged over the run, 67.4\,\% and 67.8\,\%. These differences are inside the noise of a 6-model mean (about $\pm$1.2 points). Neither run beat its start at any curve point, so both saved the round-0 champion, which is \model{greedy-full-v1} itself: the two final tests above are the same model's and therefore identical (66.7\,\% points).

\section{Behaviour}
How the climber spent its rounds (from \texttt{rounds.jsonl}):

\begin{center}\sffamily\small\begin{tabular}{@{}lrrrr@{}}\toprule
Run & Random rounds & Mirror rounds & Credit rounds & Kept\\\midrule
Mirror + credit & 58 & 47 & 210 & 2 (1 random, 1 credit)\\
Control & 318 & -- & -- & 3 (all random)\\
\bottomrule\end{tabular}\end{center}

Credit chose two thirds of the rounds and produced one keep; the 47 mirrored rounds produced none.
\todo{jakob}{What you saw in the live games, if you watched any.}

\section{Interpretation}
The bottleneck was not which knobs were tried but whether a round can see a gain at all: a round of 300--400 games cannot detect changes under about 1 point, so almost every challenger, steered or not, fails the test, and the credit model learns from margins that are mostly noise. Mirroring a clear loss assumes the knob effect is roughly linear and the loss was real; with these sample sizes most ``clear'' losses are noise as well. The same budget limit later led to sweeps per mechanism and to the perturbation probe (experiment 12), which tests every knob's z-score over many rounds instead of single challengers.
\todo{jakob}{Your reading: agree, disagree, what is missing?}

\section{Threats to validity}
One run per arm with one seed: a difference of one keep (2 vs 3) is not evidence either way. The credit prior and exploration share were not swept, so a better-tuned credit might do better. The start, \model{greedy-full-v1}, was already a 100,000-game hill-climbing result, which may leave little for any climber to find.

\section{Verdict and next step}
Not adopted: by the rule for tuning runs (the mean at 100,000 games at least the base's and the control's), mirror + credit's 66.6\,\% is below the start's 67.7\,\%, and it kept no more changes than the control. The options stay in \texttt{train-hill.ts}, off by default. The champion is unchanged. The next step taken was a different one: one mechanism per run, measured by a strength sweep (experiments 1--11), and the perturbation probe for tuning (experiment 12).

\section{Notes for the paper}
A negative result worth one line: steering a hill climber from its own past rounds does not help when single rounds cannot resolve the effect sizes involved; the limit is statistical power per round, not the proposal distribution. The table in Behaviour carries it.
\todo{jakob}{Is this worth a place in the paper? Where?}

\end{document}
